\scanpage{281}
\mspage{11}

On en \unsure{conclura|conclut} les énoncés duaux concernant les sommes directes et les objets initiaux.

La démonstration est un cas particulier du résultat bien utile suivant, appliqué au couple de foncteurs adjoints

\diagram{\begin{tikzcd}[ampersand replacement=\&,column sep=large]
M \arrow[r,shift left,"{\mathrm{diag}}"] \& M^I \arrow[l,shift left,"{\substack{\text{produits}\\\text{\unsure{dans|de} }M\\\text{indexés par }I}}" ] \& \mathpunct{:}
\end{tikzcd}}

% slanted
\margin{cf $\overline{\underline{\mathrm{XII}}}$ lemme \unsure{page} 248}

\uline{Lemme} 2.6. Soient $(M,W),(M',W')$ deux catégories de modèles, et

\diagram{\begin{tikzcd}[ampersand replacement=\&]
M \arrow[r,shift left,dash,"{\unsure{\varphi|\phi}}"] \& M' \arrow[l,shift left,"\psi"]
\end{tikzcd}}

un couple de foncteurs adjoints, compatibles aux localiseurs i.e. tels que

\[
\varphi(W)\subset W',\qquad \psi(W')\subset W.
\]

Alors les foncteurs \deleted{\unsure{induits}}\unsure{localisés}

\diagram{\begin{tikzcd}[ampersand replacement=\&]
MW^{-1} \arrow[r,shift left,"{\bar{\varphi}}"] \& M'W'^{-1} \arrow[l,shift left,"{\bar{\psi}}"]
\end{tikzcd}}

sont également adjoints, les flèches d'adjonction \inserted{pour $\bar{\varphi}\bar{\psi}$ et $\bar{\psi}\bar{\varphi}$} étant déduites de celles pour $\varphi\psi$ et $\psi\varphi$ par composition avec les foncteurs de localisation $M\unsure{{-}|\longrightarrow}MW^{-1}$, $M'{-}M'W'^{-1}$.

% archive sheet number: 280
